% Classe do documento e parâmetros gerais.
\documentclass[a4paper,openright,twoside,11pt]{report}

% Packages utilizadas e respetivos parâmetros.
\usepackage[utf8]{inputenc}
\usepackage[english]{babel}

\usepackage{lipsum} % gerador de texto
\usepackage{graphicx}
\usepackage{url}
\usepackage[Algoritmo]{algorithm}
\usepackage{algorithmicx}
\usepackage{algpseudocode}
\usepackage[printonlyused]{acronym}
\usepackage{listings}
\usepackage{hyperref}
\renewcommand{\algorithmicrequire}{\textbf{Dados: }}
\renewcommand{\algorithmicensure}{\textbf{Resultado: }}

% Definições das dimensões das páginas
\setlength{\textheight}{24.00cm}
\setlength{\textwidth}{15.50cm}
\setlength{\topmargin}{0.35cm}
\setlength{\headheight}{0cm}
\setlength{\headsep}{0cm}
\setlength{\oddsidemargin}{0.25cm}
\setlength{\evensidemargin}{0.25cm}

%\renewcommand{\baselinestretch}{1}

% Página inicial (capa)
\title{
    \vspace{-50mm}
    \begin{minipage}[l]{\textwidth}
        \hspace{-20mm}\resizebox{75mm}{!}{\includegraphics{./figures/logoISELnew2}}\\
    \end{minipage}\\[20mm]
    {\Huge \textbf{Krometheus}}\\[3mm]
    {\Large A Prometheus Kotlin client library}\\[10mm]
}

% Nome dos autores (um por linha)
\author{
    \begin{tabular}{ll}
        & Mário Carvalho \\
        & Rafael Nicolau \\[50mm]
    \end{tabular}}

\date{
    \begin{tabular}{ll}
    {Supervisor:}
        & José Simão \\
    \end{tabular}\\[10mm]
% Deixar o indicador respetivo em função da versão do relatório.
    Final report written for Project and Seminary\\
    BSc in Computer Science and Computer Engineering\\[20mm]
    June 2025}


\begin{document}
    \thispagestyle{empty}
    \maketitle

    \baselineskip 18pt % line spacing: 12pt for single, 18pt for 1 1/2, and 24pt for double spacing

    \newpage
    \thispagestyle{empty}
% Fim da contracapa

% Página com identificação completa (número e nome) e assinaturas do(s) estudante(s) e do(s) orientador(es)
    \cleardoublepage
    \setcounter{page}{1}
    \begin{center}
        \textsc{\LARGE Instituto Superior de Engenharia de Lisboa}\\[50mm]

        {\Huge \textbf{Krometheus}}\\[3mm]
        {\Large A Prometheus Kotlin client library}\\[10mm]
        \begin{tabular}{rl}
            50561 & Mário Rijo Antunes Carvalho\\[10mm]
            & \rule{75mm}{0.5pt}\\[5mm]
            50546 & Rafael Sebastião Matias Nicolau\\[10mm]
            & \rule{75mm}{0.5pt}              \\
        \end{tabular}\\[10mm]

        \begin{tabular}{rl}
            Supervisor: & José Manuel de Campos Lages Garcia Simão\\[10mm]
            & \rule{75mm}{0.5pt}\\[5mm]
        \end{tabular}\\[10mm]

        Final report written for Project and Seminary\\
        BSc in Computer Science and Computer Engineering\\[20mm]
        June 2025\\
    \end{center}

% Página de resumo em ingles
    \cleardoublepage
    \chapter*{Abstract}
    The growing complexity of modern software systems, driven by the adoption of distributed architectures,
    microservices, and asynchronous execution—has made observability a critical aspect of system design.
    Among its key components, metrics-based monitoring enables developers and operators to track the internal state and performance of applications and infrastructure.
    Prometheus has become a widely used standard for collecting and querying metrics; however, its official client libraries primarily target languages such as Go, Python, and Java.
    These implementations do not offer idiomatic support for Kotlin, its coroutine-based concurrency model, or
    Kotlin-specific frameworks like Ktor.\\

    This project presents a Prometheus client library implemented in Kotlin, designed with coroutine safety, expressiveness, and alignment with Kotlin's language features in mind.
    The library supports core Prometheus metric types and provides a straightforward and testable interface for instrumenting Kotlin applications.
    This report details the motivation for the project, its design and architecture, the implementation process, and an evaluation of the resulting system.
    The library fills an existing gap in the Kotlin ecosystem by offering a native and idiomatic solution for
    integrating metrics-based observability.\\

    {\textbf{Keywords}: Observability, Prometheus, Metrics, Monitoring.}

%% Página de agradecimentos
    \cleardoublepage
    \chapter*{Acknowledgments}
    \lipsum[1]
    \vspace{2em}
    \begin{flushright}
        Rafael Nicolau
    \end{flushright}

    \lipsum[1]
    \vspace{2em}
    \begin{flushright}
        Mário Carvalho
    \end{flushright}

% Geração do índice de conteúdos
    \cleardoublepage
    \tableofcontents \cleardoublepage

% Geração do índice de figuras e de tabelas
    \listoffigures \cleardoublepage
    \listoftables \cleardoublepage
    \chapter*{List of Abbreviations}
\addcontentsline{toc}{chapter}{List of Abbreviations}

\begin{acronym}[TDMA]  % widest acronym to set the width for the list
    \acro{ARPANET}{Advanced Research Projects Agency Network}
    \acro{ENIAC}{Electronic Numerical Integrator and Computer}
    \acro{WWW}{World Wide Web}
    \acro{TCP/IP}{Transmission Control Protocol/Internet Protocol}
    \acro{FTP}{File Transfer Protocol}
    \acro{SSH}{Secure Shell}
    \acro{SOA}{Service-oriented Architecture}
    \acro{SaaS}{Software as a Service}
    \acro{EC2}{Elastic Compute Cloud}
    \acro{VM}{Virtual Machine}
    \acro{SLA}{Service Level Agreement}
    \acro{API}{Application Programming Interface}
    \acro{KMP}{Kotlin Multiplatform}
    \acro{MTTR}{Mean Time To Repair}
    \acro{CNCF}{Cloud Native Computing Foundation}
    \acro{CPU}{Central Processing Unit}
    \acro{JVM}{Java Virtual Machine}
    \acro{DSL}{Domain-Specific Language}
\end{acronym}
 \cleardoublepage

% Chapter 1
    %
% Chapter 1
%
\chapter{Introduction} \label{ch:introduction}


\section{Context and Motivation}\label{sec:context-and-motivation}
Over the past few decades, we have been witnessing remarkable advancements on our software systems
and their architecture.\\

The first computers, mainframe computers, Harvard Mark I and the \ac{ENIAC} showed up in the 1930s-1940s and where
developed for military and research purposes.
These were large, powerful hardware machines that took up to an entire room.
The software architecture for these applications was monolithic, meaning that they were simple, self-contained and
independent of other applications~\cite{orkes_software_architecture_evolution, wikipedia_monolithic}.\\

Networks connect and facilitate communication between computers--mainframe to terminal, mainframe to mainframe and
later client to server.
The development of network technology from 1958 onwards enabled mainframes to be connected
electronically, transforming them from isolated machines to multi-user computers that were connected to multiple
terminals.
\ac{ARPANET} was the first public, wide-area computer network, going live in 1969.
It communicated using
packet switching, which went on to serve as the foundation for modern-day Internet as we know it.
Network technology popularized the client-server structure in the 80s, where an application is divided into a
server and a client communicate over a network.
This is a very familiar structure to us today: a client, typically a
desktop computer, remotely makes a request to a server, which returns a response.
This way we can partition tasks
and workloads, the server will work on the processing and retrieval of a piece of data, and the client will in turn
present it~\cite{orkes_software_architecture_evolution, wikipedia_client_server_model}.\\

1983 marked the year of the Internet, a global system of computer networks that used the \ac{TCP/IP} protocol to
facilitate communication between devices and applications.
This was the backbone for \ac{FTP} programs, \ac{SSH} systems, and of course, the World Wide Web.
The invention of the web and its latent possibilities soon kicked off the next wave of application development.
Instead of building a dedicated client for your client for your application, you could simply build a website to be
hosted on the web~\cite{orkes_software_architecture_evolution}.\\

As application development grew, a monolithic codebase became more unwieldy to manage, and it became clear that
capabilities and data house in a system could be reused.
To address this pain point, modularization became a topic of discussion, and in the 90s the first N-tiered
applications showed up, for the first time we saw the server being split into two tiers: the application server and
the database, the application held all the application and business logic, while the database server, stored the
data recorded, which reduced latency at high processing volumes.
Around the same time, \ac{SOA} emerged as an architectural pattern, and we saw applications split into multiple
services, each design to perform a specific job which could be things like processing payments or verifying user
identities, this had benefits such as an increased scalability as single services could be scaled up or down
depending on demand without having to make changes to the entire system
~\cite{orkes_software_architecture_evolution, oracle_soa, port_soa}.\\

\ac{SOA} set the stage for the move from traditional desktop applications to a new mode of software applications - \ac{SaaS},
but it was the invention of virtual machines and cloud computing that further spurred the explosion of \ac{SaaS}
products in the coming decades.
Machine virtualization exists since the 1960s, but this technology only came into
mainstream use in the 2000s.
It at this period that Amazon, followed by companies like Google, Microsoft and Oracle identified the lucrative that
virtualization offered: managed cloud computing services.
With cloud computing services like Amazon \ac{EC2}, companies could rent \ac{VM}s for processing power and scale as
needed~\cite{orkes_software_architecture_evolution}.\\

The 2010s were the culmination of multiple trends towards distributed computing.
Fueled by the need for third-party access to their services, the first commercial \ac{API}s were lauched in the
2000 by Salesforce and eBay, to enable their costumers and partners to integrate features onto their own sites and
applications.
From Google Maps and Stripe to Twilio and OpenAI, the \ac{API} economy has ballooned since, powering integrated
features across the web.
In the same vein, microservices took off when scaling companies like Netflix and Amazon needed to speed up and
streamline the development cycle, which was slowed by a monolithic architecture.
By splitting up an application into individual microservices, each with its own database, teams could independently update and deploy them, leading to
faster releases and improvements.\cite{orkes_software_architecture_evolution}
The growth in the adoption of this architectural style was heavily supported by the emergence of containers.\\

As we move into the present decade, software systems have become increasingly complex and distributed.
Modern applications often consist of hundreds or thousands of microservices, deployed across hybrid and multi-cloud environments, communicating asynchronously via APIs, message queues, and event streams.
Containers and orchestration platforms like Kubernetes have become the de facto standard for managing these
distributed components at scale\cite{devtron_kubernetes}.
While this architectural evolution brings immense benefits—such as scalability, resilience, and faster development cycles—it also introduces significant operational challenges.
The large number of independent moving parts, dynamic scaling, and network dependencies make it difficult to predict
and diagnose failures or performance degradation using traditional reactive methods.\\

Consequently, couple with more stringent \ac{SLA}s than ever, for multiple reasons, organizations now emphasize
observability, the
ability to
measure a
system's
current
state based
on its
external outputs, recorded as metrics, logs, and traces\cite{dynatrace_observability}.
Observability enables a proactive approach to maintaining a systems reliability and performance standards, allowing
teams to detect
anomalies
early,
understand complex interactions, and respond swiftly before issues impact end users.\\


In this context, metrics-based monitoring, has become essential for maintaining the health of modern distributed systems.
This project aims to fill a notable gap in the Kotlin ecosystem by providing a robust and idiomatic Prometheus client library tailored specifically for Kotlin applications.
Unlike existing libraries that primarily target other JVM languages or lack native support for Kotlin’s coroutine model, this library offers seamless integration and out-of-the-box support for popular Kotlin-only frameworks such as Ktor.
By delivering a solution designed with Kotlin’s language features and concurrency paradigms in mind, this project enables developers to instrument their applications effectively and embrace modern metrics-based monitoring with minimal friction.

\section{Objectives}\label{sec:objectives}

The primary objective of this project is to develop a Kotlin-native Prometheus client library that addresses existing limitations in the Kotlin ecosystem by providing a robust, idiomatic, and efficient solution for metrics instrumentation and collection.

Specifically, the project aims to:

\begin{itemize}
    \item Design and implement a client library that fully leverages Kotlin’s coroutine-based concurrency model to enable asynchronous and efficient metric collection.
    \item Provide seamless integration with popular Kotlin frameworks such as \texttt{Ktor}, ensuring ease of use and
    compatibility within typical Kotlin application environments.
    \item Ensure the library is designed with readability and idiomatic Kotlin style in mind, making it intuitive and enjoyable to use.
    \item Develop an expressive and idiomatic \ac{DSL} for defining and managing metrics,
    reducing boilerplate code and improving developer productivity.
    \item Optimize the library for performance and resource efficiency, surpassing existing Java-based Prometheus
    clients when used in Kotlin applications where possible.
    \item Support \ac{KMP}, enabling broad platform compatibility and facilitating code reuse across \ac{JVM},
    Android, and
    other supported targets.
\end{itemize}

By achieving these objectives, the project intends to fill a critical gap in Kotlin observability tooling and empower developers to instrument their applications with minimal effort while adhering to modern Kotlin programming paradigms.

\section{Structure of the Report}\label{sec:structure-of-the-report}

This report is organized into eight main chapters, each focusing on a specific aspect of the project:

\begin{itemize}
    \item \textbf{Chapter 1: Introduction} — Presents the context and motivation behind the project, outlines the objectives, and provides an overview of the report’s structure.

    \item \textbf{Chapter 2: Background and Motivation} — Covers foundational concepts related to observability and Prometheus, reviews existing Prometheus client libraries, and highlights gaps within the Kotlin ecosystem that motivate this work.

    \item \textbf{Chapter 3: Problem Statement and Goals} — Defines the specific limitations identified in current tooling and states the precise goals and scope of the project.

    \item \textbf{Chapter 4: Design and Architecture} — Describes the high-level architectural decisions, focusing on concurrency models, coroutine usage, and the lifecycle of metric registries.

    \item \textbf{Chapter 5: Implementation Details} — Details the concrete implementation of the client library, including metric types, collector logic, and framework integrations such as with Ktor.

    \item \textbf{Chapter 6: Evaluation and Limitations} — Presents the functional validation, performance evaluation, and comparison with other clients, as well as known limitations of the project.

    \item \textbf{Chapter 7: Future Work} — Suggests possible extensions and improvements, including support for additional technologies, enhanced integration with Kotlin frameworks, and automatic JVM metrics collection.

    \item \textbf{Chapter 8: Conclusion} — Summarizes the contributions made by the project, reflects on lessons learned, and offers final remarks.
\end{itemize}

The report concludes with a list of references and an appendix containing supplementary material.


% Chapter 2
    %
% Chapter 2
%
\chapter{Background and Motivation} \label{ch:background}


\section{Observability Overview}
\lipsum[1]


\section{Prometheus Overview}
\lipsum[1]


\section{Existing Prometheus Clients}
\lipsum[1]


\section{The Kotlin Ecosystem and Observability Gaps}
\lipsum[1]

% Chapter 3
    \chapter{Problem Statement and Goals} \label{ch:problemdescription}

\section{Identified Limitations}

While Prometheus has become the de facto standard for metrics-based observability in cloud-native systems, its official client libraries are largely written in imperative languages such as Go, Java, and Python. These implementations do not align well with Kotlin’s idiomatic programming model, particularly regarding coroutine support, functional DSLs, and thread-safety via structured concurrency.

The official Java Prometheus client, while usable in Kotlin projects, exposes a verbose and boilerplate-heavy API. Furthermore, it does not integrate naturally with Kotlin coroutines, leading to potential blocking behavior or unsafe metric updates when used improperly in suspendable contexts.

Additionally, Kotlin developers lack a native library that fully embraces Kotlin-specific features like extension functions, scope builders, and coroutine-friendly primitives for metrics definition and recording. This mismatch creates friction when building observability into modern reactive applications using frameworks such as Ktor or Spring WebFlux.

\section{Project Goals}

This project aims to develop a lightweight, idiomatic Prometheus client library written in Kotlin. The primary goals are:

\begin{itemize}
  \item \textbf{Kotlin-first design:} Create a metrics API that aligns with Kotlin’s language features, using builder functions, lambdas, and extension patterns to reduce boilerplate and improve readability.
  \item \textbf{Coroutine safety:} Ensure that metric operations can be safely invoked from within suspending functions and concurrent coroutines, without introducing race conditions or blocking behavior.
  \item \textbf{OpenMetrics compatibility:} Conform to the OpenMetrics exposition format, ensuring compatibility with Prometheus scraping, tooling, and external monitoring pipelines.
  \item \textbf{Core metric support:} Implement the most commonly used metric types—\texttt{counter}, \texttt{gauge}, \texttt{histogram}, and \texttt{summary}—with support for labeling, concurrency, and serialization.
  \item \textbf{Extensibility:} Provide a modular and extensible architecture, allowing future features such as HTTP exposition, JVM metrics integration, or PushGateway support to be added incrementally.
\end{itemize}

\section{Out of Scope}

To maintain focus and deliver a stable MVP, certain features and use cases are intentionally excluded from this project:

\begin{itemize}
  \item \textbf{PushGateway integration:} The library will support pull-based scraping only. Push-based metric transport (e.g., batch job metrics) is not addressed.
  \item \textbf{Advanced metric types:} Specialized types like \texttt{info}, \texttt{enum}, or \texttt{state\_set} will not be implemented in the initial version.
  \item \textbf{JVM internals instrumentation:} Metrics related to garbage collection, memory pools, or thread activity are considered outside the scope of this project.
\end{itemize}


% Chapter 4
    \chapter{Design and Architecture} \label{ch:architecture}

\section{High-Level Architecture}

The Prometheus client for Kotlin is designed with a modular, coroutine-friendly architecture that adheres to the \textbf{OpenMetrics specification}~\cite{openmetrics-spec}. This design enables robust metric definition, collection, and exposition in distributed Kotlin applications while remaining compatible with modern observability tooling such as Prometheus.

At its core, the architecture separates responsibilities into clearly defined layers:

\begin{itemize}
    \item \textbf{Collector Abstractions:} The base interfaces and abstract classes that define how metrics are described and collected.
    
    \item \textbf{Concrete Metric Types:} Implementations such as \texttt{Counter} inherit from a shared label-aware base class and enforce the conventions required by OpenMetrics (e.g., \texttt{\_total} suffix, monotonicity for counters).
    
    \item \textbf{Label Management:} A child-metric system that maps unique sets of label values to individual metric instances.
    
    \item \textbf{Registry:} A coroutine-safe global registry that allows metrics to be registered, queried, or cleared. 
    
    \item \textbf{Text Formatting Layer:} A dedicated formatter that converts collected metrics into the OpenMetrics text exposition format. This includes headers such as \texttt{\# TYPE}, \texttt{\# HELP}, and \texttt{\# UNIT}, and optionally includes high-resolution timestamps as specified by the standard.
    
    \item \textbf{Exporter:} A simple front-facing API that aggregates all registered metrics and returns a complete OpenMetrics-compliant snapshot for scraping.
\end{itemize}

\section{Concurrency and Coroutines}

The Prometheus client for Kotlin is designed to operate efficiently and safely in concurrent environments common to modern Kotlin applications. To achieve this, it leverages Kotlin's coroutines and atomic primitives from the \texttt{kotlinx.atomicfu} library to ensure lock-free, thread-safe updates to metric values.

\subsection{Atomic Metric Updates}

Metrics such as counters must be safely incremented across multiple threads or coroutines without losing updates or causing data races. Instead of using traditional locks or mutexes, the client employs atomic operations on the raw bits of \texttt{Double} values, providing a performant non-blocking solution.

For example, the \texttt{Counter} class internally stores its value as an \texttt{atomic} long representing the bit pattern of a \texttt{Double}. Increment operations update this atomic value using \texttt{updateAndGet} with a lambda that atomically reads the current bits, converts them to a \texttt{Double}, adds the increment, and writes back the updated bits.

\subsection{Coroutine Context Offloading}

To prevent blocking critical coroutine dispatchers or event loops, the increment operation is dispatched onto \texttt{Dispatchers.Default} using \texttt{withContext}. This design choice balances safety and performance by offloading potentially compute-heavy atomic updates to a background thread pool optimized for CPU-bound tasks.

\subsection{Suspendable Increments}

Increment functions such as \texttt{inc(amount: Double)} are declared \texttt{suspend} to integrate seamlessly with asynchronous code. This allows calling coroutines to perform metric updates without blocking the main thread or requiring additional synchronization primitives.

\subsection{Generalization to Other Metric Types}

While this concurrency approach is illustrated with the \texttt{Counter} metric, it applies equally to other metric types like \texttt{Gauge} that require thread-safe mutation. Each child metric instance independently manages its atomic state, enabling concurrent, lock-free updates even in heavily parallelized applications.

\subsection{Summary}

By combining atomic primitives and coroutine dispatchers, the Kotlin Prometheus client achieves a concurrency model that is both safe and performant. This design enables reliable metric collection in asynchronous, multi-threaded environments without sacrificing responsiveness or introducing complex locking logic.
\section{Metric Registry and Lifecycle}

The lifecycle of a metric in the Kotlin Prometheus client spans from its initial creation to the eventual collection of its samples during scraping. Understanding this lifecycle clarifies how metrics are instantiated, managed, and exposed within the system.

\subsection{Metric Creation}

Metrics are created by instantiating concrete subclasses of the abstract \texttt{Collector} class, such as \texttt{Counter} or \texttt{Gauge}. Each metric requires key metadata including a fully qualified name, a descriptive help string, a list of label names, and optionally a unit of measurement.

Upon creation, metrics validate their names and label names against OpenMetrics-compliant regular expressions to ensure compatibility with monitoring systems.

\subsection{SimpleCollector and Child Metrics}

The client leverages a generic \texttt{SimpleCollector} abstraction that manages label-based child metrics. Each unique combination of label values corresponds to a \texttt{Child} instance that holds the actual metric state (e.g., a counter’s value).

When a metric is accessed with specific label values via the \texttt{labels()} method, the collector either retrieves an existing \texttt{Child} or creates a new one. This lazy initialization ensures efficient memory usage by only instantiating metric instances as needed.

\subsection{Registry Registration}

Once created, metrics are registered with a \texttt{CollectorRegistry}, either explicitly or via a default global registry. Registration guarantees the uniqueness of metric names within the registry and enables centralized management.

The registry maintains a thread-safe collection of all registered metrics, using a coroutine-friendly \texttt{Mutex} to synchronize additions and removals. This approach supports dynamic lifecycle management, allowing metrics to be registered or unregistered during runtime.

\subsection{Metric Mutation and Updates}

Metric values are updated through their respective \texttt{Child} instances, which expose safe, concurrent operations such as incrementing counters or setting gauges. Atomic operations and coroutine context switching ensure thread-safe, lock-free updates even in asynchronous environments.

\subsection{Sample Creation and Collection}

During a scrape, the registry invokes the \texttt{collect()} method on each registered metric. This triggers the metric to aggregate its current state into a list of \texttt{Sample} instances, each representing a single data point with its name, labels, value, and optional timestamp.

These samples are gathered into a \texttt{MetricFamilySamples} object, which encapsulates the entire metric family with metadata such as type and help text, in accordance with the OpenMetrics specification.

\subsection{Exposition}

Finally, the collected samples are passed to a formatter that converts them into the OpenMetrics text exposition format. This formatted output is then served to monitoring systems for ingestion and visualization.

\bigskip

Through this lifecycle—from creation, label-based child instantiation, registry management, concurrent updates, to sample aggregation—the client ensures accurate, consistent, and efficient metric reporting that integrates seamlessly with Prometheus and OpenMetrics-compatible tools.


% Chapter 5
    \chapter{Implementation Details} \label{cap:implementation}


\section{Metric Types and API Design}
\lipsum[1]


\section{Collector Registration and Logic}
\lipsum[1]


\section{Ktor Integration}
\lipsum[1]


% Chapter 6
    \chapter{Evaluation and Limitations} \label{cap:evaluation}

\section{Functional Validation}

The Kotlin Prometheus client was functionally validated against the \textit{OpenMetrics} specification~\cite{openmetrics-spec}, ensuring compliance with required formatting rules, metadata emission, and semantic conventions.

To validate correctness, unit tests were written to verify:
\begin{itemize}
  \item Label handling and metric uniqueness.
  \item Output formatting according to OpenMetrics text exposition format.
  \item Correct behavior of metric operations (\texttt{inc}, \texttt{get}, \texttt{set}) across gauges, counters, histograms, and summaries.
  \item Accurate quantile configuration and observation bucketing for summaries and histograms.
  \item Safe concurrent access using coroutines and \texttt{atomicfu} primitives.
\end{itemize}

Support for histograms includes both linear and exponential bucket configuration. Summary metrics were validated for configurable quantile tracking and tolerances. The client was successfully scraped using a Prometheus server, confirming its compatibility with external collectors and exporters.

\section{Performance Considerations}

The Kotlin client prioritizes coroutine-friendly concurrency and safe metric mutation over raw throughput. As shown in benchmarking (see Chapter~\ref{ch:architecture}), read-heavy operations (\texttt{get}) perform exceptionally well and sometimes surpass the Java client in throughput.

Write-heavy operations such as \texttt{inc()}, \texttt{set()}, and \texttt{observe()} demonstrate slightly lower performance due to coroutine scheduling and atomic synchronization overhead. Histograms and summaries involve more complex logic (e.g., bucketing, rank approximation), which adds some cost to observation recording. However, the performance remains acceptable and predictable under load.

Despite this overhead, the client maintains stable performance under high concurrency thanks to:
\begin{itemize}
  \item Non-blocking metric access patterns.
  \item Scoped coroutine isolation for metric operations.
  \item Use of lock-free atomic primitives.
\end{itemize}

\section{Comparison with Other Clients}

When compared to the official Java Prometheus client:
\begin{itemize}
  \item The Kotlin client has a more idiomatic API tailored to Kotlin, with extension functions, lambdas, and DSL-style builders.
  \item Coroutine support enables suspendable metric operations and better alignment with modern Kotlin concurrency models.
  \item Histograms and summaries follow a simpler configuration model while remaining fully compliant with OpenMetrics.
  \item Mutation operations are slightly slower due to Kotlin's atomic wrappers instead of Java's \texttt{LongAdder}.
\end{itemize}

The inclusion of histogram and summary support brings metric coverage on par with the Java client, while preserving the Kotlin-first design philosophy. The implementation was validated against expected output structures and snapshot comparisons with Java-based benchmarks.

\section{Known Limitations}

While the Kotlin Prometheus client provides robust core features and coroutine support, some limitations remain:
\begin{itemize}
  \item \textbf{Limited metric types:} The client currently supports only core metric types—\texttt{counter}, \texttt{gauge}, \texttt{histogram}, and \texttt{summary}. More specialized metrics such as \texttt{info}, \texttt{state\_set}, or \texttt{enum} are not implemented.
  \item \textbf{No use of \texttt{LongAdder}:} The client avoids using JVM-specific classes like \texttt{LongAdder}, which offer superior write throughput in high-concurrency environments. Instead, it relies on portable atomic primitives to ensure coroutine safety, at the cost of some raw performance.
  \item \textbf{No PushGateway support:} Push-based metric delivery to Prometheus (e.g., via PushGateway) is not currently supported.
\end{itemize}

Future enhancements could address these limitations by:
\begin{itemize}
  \item Adding support for PushGateway integration.
  \item Extending metric coverage to include specialized types such as \texttt{info} and \texttt{state\_set}.
  \item Exploring conditional use of \texttt{LongAdder} via Java interop for high-performance counters.  \item Supporting dynamic histogram configuration.
\end{itemize}



% Chapter 7
    \chapter{Future Work} \label{cap:future}

\section{Expanding Ecosystem Integrations}
While the current Prometheus client supports integration with Ktor, future development should focus on extending support to other popular server frameworks and platforms within the Kotlin ecosystem. This includes native support for:
\begin{itemize}
    \item \textbf{Spring Boot (Kotlin DSL):} Enabling Prometheus metrics registration via auto-configuration and Spring Actuator endpoints.
    \item \textbf{Micronaut and Koin:} Providing lightweight instrumentation tailored for dependency-injection-heavy applications.
    \item \textbf{Android and Kotlin Multiplatform:} Extending observability tooling to mobile or cross-platform environments where Prometheus scraping is less common, potentially using push-based collection or local aggregators.
\end{itemize}

Such integrations would promote broader adoption by reducing the friction for developers building on top of different stacks.

\section{Improve Integration with Well-Known Kotlin Frameworks}
The current Ktor integration demonstrates the feasibility of lightweight, coroutine-friendly metric collection. However, more ergonomic integrations can be achieved by:
\begin{itemize}
    \item Introducing \textbf{Kotlin-style DSLs} for declaratively registering metrics at application startup.
    \item Offering built-in support for Ktor’s typed routing, allowing route-specific metrics without requiring manual normalization or labeling.
    \item Enabling pluggable interceptors to monitor application features like database access, authentication, or external API calls.
\end{itemize}

Providing idiomatic APIs that reflect Kotlin’s language features (extension functions, delegation, lambdas) would help developers embed metrics more naturally into their applications.

\section{Automatic JVM Metrics Collection}
A key enhancement for production-readiness is enabling automatic collection of runtime-level metrics from the JVM environment. This includes:
\begin{itemize}
    \item \textbf{Garbage collection and memory usage:} Exposing heap/non-heap statistics and GC pause durations.
    \item \textbf{Thread metrics:} Monitoring live threads, daemon threads, and blocked thread counts.
    \item \textbf{Class loading and compilation metrics:} Tracking class definitions, unloads, and JIT compilation time.
\end{itemize}

These metrics are critical for performance tuning and operational visibility, especially in cloud-native deployments. Leveraging the standard \texttt{java.lang.management} API, these metrics can be collected automatically on JVM startup and registered in the Prometheus client’s registry.

Combining application-level and system-level metrics would provide a complete observability picture, aligning the client with the capabilities offered by the official Prometheus Java client.


% Chapter 8
    \usepackage{lipsum}


\chapter{Conclusion} \label{cap:conclusion}


\section{Summary of Contributions}
\lipsum[1]


\section{Lessons Learned}
\lipsum[1]


\section{Final Remarks}
\lipsum[1]

% References
    \bibliographystyle{plain}
    \bibliography{refs}
    \addcontentsline{toc}{chapter}{References}

    \appendix
    %
% Appendix 1
%
\chapter{Appendix} \label{ap:exemplo}
\lipsum[14-16]

\end{document} 